\section{Weak asymptotics in four dimensions}
\label{n4:sec:introduction}\label{sec:weak4}

We prove a lower bound for the ADM rest mass in terms of the least
three-volume needed to separate a chosen asymptotically flat end from
the compact weakly trapped boundary and the other ends. The task is to
bring two derivatives of the metric and one derivative
of the second fundamental form within the scope of
Theorem~\ref{thm:numerical:prepared}. We first replace the selected end
by a rest end while controlling every cut, then make the dominant
energy and boundary inequalities strict. The prepared theorem is
applied to each fixed resulting datum; only its numerical conclusion
is passed to the limit.

\subsection{Initial data and enclosing cuts}

Four is the spatial dimension; the associated spacetime dimension is
five, and a hypersurface area is a three-dimensional Riemannian volume.
Here $\omega_3=|S^3|=2\pi^2$, and we write
$\tau=\operatorname{tr}_gK$. Specializing the constraint convention
already fixed in the paper gives
\[
 2\mu=R_g+(\operatorname{tr}_gK)^2-|K|_g^2,\qquad
 J_i=\nabla^j(K_{ij}-(\operatorname{tr}_gK)g_{ij}).
\]
Here $R_g=\Scal_g$ denotes scalar curvature and $\nabla$ is the
Levi--Civita connection of $g$. The cosmological constant is zero. We use
\(\Delta_g=\divg_g\grad_g\).
For a hypersurface with specified unit normal \(\nu\), write
\[
 H=\divg_{\mathrm{tan}}\nu,\qquad
 \tr_{\mathrm{tan}}K=(g^{ij}-\nu^i\nu^j)K_{ij},\qquad
 \Theta_K=H+\tr_{\mathrm{tan}}K.
\]
At an inner boundary the normal points into the exterior. In an auxiliary
trapped-region construction it points out of the bounded region.

The notation \(O_j(r^{-a})\) includes coordinate derivatives of order
\(k\le j\), bounded by \(O(r^{-a-k})\). On an asymptotically flat end the
ADM normalization is
\begin{align}
 E&=\frac{1}{6\omega_3}\lim_{R\to\infty}
       \int_{|x|=R}(\partial_jg_{ij}-\partial_i g_{jj})
                    n_\delta^i\,\dd A_\delta,\label{n4:intro:energy}\\
 P_i&=\frac{1}{3\omega_3}\lim_{R\to\infty}
       \int_{|x|=R}(K_{ij}-\tau g_{ij})
                    n_\delta^j\,\dd A_\delta.\label{n4:intro:momentum}
\end{align}
Here \(n_\delta\) and \(\dd A_\delta\) are Euclidean.
Whenever a charge is used below, its stated limit is required to exist.

\begin{definition}[Enclosing cuts]\label{n4:def:cuts}
Let \(M\) have nonempty compact boundary \(S\) and finitely many
asymptotically flat ends, and fix an end \(e\).
An admissible outer domain \(D_e\) is a connected smooth
codimension-zero submanifold with boundary, closed as a subset of \(M\),
such that its manifold interior lies in \(\operatorname{int}M\),
it contains the entire sufficiently distant part of \(e\), and
it contains no sufficiently distant part of any other end.
Its entire intrinsic manifold boundary
\(\Gamma=\partial_{\mathrm{man}}D_e\) is required to be compact.
The enclosing volume is
\[
 A_e(S)=\inf_{D_e}|\partial_{\mathrm{man}}D_e|_g.
\]
All frontier coincident with \(S\) is counted. The cut may be
disconnected. In the one-ended case \(D_e=M\) is allowed and has
intrinsic boundary \(S\). We then also write \(A_*\) for this infimum.
\end{definition}

Only these smooth cuts are used. Their infimum need not be attained.
The connected outer domain excludes the distant parts of all other
ends, although its cut may be disconnected. Thus all other ends and
the entire original boundary are obstacles
to the selected outer domain. This end-isolating convention remains
distinct from the three-dimensional complete-data class, whose
chosen-end region may be disconnected and may retain other ends.

\begin{theorem}\label{n4:thm:weak}
Let \((M^4,g,K)\) be smooth and
connected, with \(M\) orientable, \(K\) an arbitrary smooth symmetric
covariant two-tensor, and \(S=\partial M\) nonempty, compact, and smooth.
Suppose \(g\) is complete as a metric space with \(S\) included.
Outside a compact set let \(M\) have finitely many, and at least one,
ends diffeomorphic to the complement of a closed ball in \(\R^4\).
Assume
\[
 \mu\ge|J|_g,\qquad \mu,|J|_g\in L^1(M,\dd V_g),
\]
and, on every end for a common \(q>1\),
\[
 g-\delta=O_2(r^{-q}),\qquad K=O_1(r^{-1-q}),
\]
with finite energy and momentum limits
\eqref{n4:intro:energy}--\eqref{n4:intro:momentum}.
On every component of \(S\) assume \(\Theta_K\le0\), with normal into \(M\).

Then for each end \(e\),
\begin{equation}\label{n4:intro:main}
 E_e\ge
 \sqrt{|P_e|_\delta^2+
       \frac14\left(\frac{A_e(S)}{\omega_3}\right)^{4/3}}.
\end{equation}
In particular \(E_e>|P_e|_\delta\), and the invariant mass satisfies
\begin{equation}\label{n4:intro:mass}
 \sqrt{E_e^2-|P_e|_\delta^2}
   \ge\frac12\left(\frac{A_e(S)}{\omega_3}\right)^{2/3}.
\end{equation}
\end{theorem}

The strict timelikeness is a conclusion: Lemma~\ref{n4:lem:cut-reduction}
establishes \(A_e(S)>0\), and the proof includes the non-timelike case.
There is no extra decay assumption on \(\tau=\tr_gK\).
All original boundary components use the same future trapping convention.

The coefficient is sharp. On the time-symmetric
Schwarzschild--Tangherlini exterior \cite{Tangherlini1963},
\[
 g_m=\left(1+\frac{m}{2r^2}\right)^2\delta,\qquad
 K=0,\qquad r\ge\sqrt{m/2},
\]
the ADM energy is \(m\), the momentum is zero, and the boundary
has three-volume \(\omega_3(2m)^{3/2}\).
The areal radius \(r+m/(2r)\) is nondecreasing on this exterior.
Radial projection to its boundary is therefore nonexpanding on
three-dimensional tangent volumes; its degree on an enclosing cut is one.
The minimum enclosing volume is consequently the boundary volume,
and \eqref{n4:intro:main} is an equality for this family.
The theorem concerns the numerical bound. Classification at equality
requires a separate argument on the original data.

\subsection{The exact cut class and removal of other ends}

\begin{lemma}[Cut reduction and positivity]\label{n4:lem:cut-reduction}
Under the hypotheses of Theorem~\ref{n4:thm:weak}, the quantity \(A_e(S)\)
is finite and strictly positive for every end \(e\).
Truncating all other ends at sufficiently large coordinate spheres gives
a connected, complete one-ended exterior \(M_T\) with compact smooth
boundary \(S_T\). All its boundary components have future nonpositive
expansion with the normal into \(M_T\). Its selected-end charges are the
original charges, and
\begin{equation}\label{n4:end:cut-truncation}
 A_*(M_T,g)\ge A_e(S),\qquad
 \lim_{T\to\infty}A_*(M_T,g)=A_e(S).
\end{equation}
The common parameter \(T\) in the limit can be replaced by any exhaustion
in which every unwanted truncation radius tends to infinity.
\end{lemma}

\begin{proof}
A sufficiently distant sphere in the chosen end bounds an admissible
outer domain, so the infimum is finite. To obtain a positive lower bound,
choose an embedded arc from an open disk in \(S\) to the chosen end, and
continue it along a straight coordinate ray. The compact part of the arc
can be chosen embedded and disjoint from the boundary except at its
initial point. A boundary collar and the tubular neighborhood theorem
give a proper tube with coordinates
\[
 [0,\infty)\times B^3\longrightarrow M.
\]
Its initial section lies in \(S\); its far part is a Euclidean tube of
fixed transverse coordinate size. Choose a smooth nonnegative function
\(\eta\) compactly supported in \(B^3\), with integral one, and define
in this tube
\[
 \alpha=\eta(z)\,\dd z^1\wedge\dd z^2\wedge\dd z^3.
\]
Extend \(\alpha\) by zero across the lateral tube boundary. It is a
smooth closed three-form on \(M\). Its comass, meaning the supremum of
its absolute value on unit simple three-vectors, has a finite bound
\(C_\alpha\): the initial part is compact, and the far metric is
uniformly comparable with the Euclidean metric. A sufficiently distant
selected-end sphere has \(\alpha\)-integral one, after fixing orientation.

Let \(D_e\) be any admissible outer domain with cut \(\Gamma\). Cut off
its selected-end tail at a sphere beyond \(\Gamma\). The resulting
domain is compact, since \(D_e\) contains no distant part of another
end. Its boundary is the distant sphere together with the \emph{entire}
intrinsic boundary \(\Gamma\), with the boundary orientations.
Stokes' theorem gives
\[
 \left|\int_\Gamma\alpha\right|=1,
 \qquad |\Gamma|_g\ge C_\alpha^{-1}.
\]
This proof applies to disconnected cuts and to frontier on \(S\).
In particular, it proves positivity before any end is truncated.

On an unwanted end, the normal of the truncation sphere that points
into the retained manifold is the inward radial normal. The given decay
therefore gives
\[
 H=-3T^{-1}+O(T^{-1-q}),\qquad
 \tr_{S_T}K=O(T^{-1-q}).
\]
The new boundary has strictly negative future expansion for large
\(T\). The other boundary components are the original ones.
Removing the open tails leaves a connected manifold: a path entering a
removed tail can have that portion replaced by a path along its connected
spherical cross-section. The retained manifold is closed in \(M\).
It is complete for its intrinsic metric as well. Indeed, an intrinsic
Cauchy sequence is ambient Cauchy and has a limit in the retained set;
near a smooth boundary or an interior point, the intrinsic and ambient
local distances are comparable in a half-ball or ball chart.

Every admissible outer domain in \(M_T\), viewed as a submanifold of
\(M\), is an admissible outer domain for \(e\). A part of its frontier
on a truncation sphere becomes a hypersurface in the interior of \(M\),
and remains part of its intrinsic boundary. This proves the first
inequality in \eqref{n4:end:cut-truncation}. Conversely, a fixed original
admissible domain has compact frontier and contains no tail of any
unwanted end. It is consequently contained in \(M_T\) for all
sufficiently large truncation radii and is then an admissible domain
there. Applying this observation to cuts with area within any prescribed
positive error of \(A_e(S)\) proves the limit. The selected-end data have
not been changed, so neither selected-end charge changes.
\end{proof}

We may thus work on one-ended data until the final transfer. The boundary
is always included. No finite-perimeter relaxation or minimizing cut is
needed in this section.

\subsection{Conformal changes and an energy-raising shell}

\begin{lemma}[Conformal constraint and boundary formulas]\label{n4:lem:conformal}
Let \(u>0\) be smooth, and set \(g_u=u^2g\), \(K_u=uK\). In dimension
four, as an identity of scalar functions and an identity of covectors,
respectively,
\begin{align}
 u^2\mu_u&=\mu-3u^{-1}\Delta_g u,\label{n4:end:conformal-mu}\\
 uJ_u&=J+3K(\grad_g\log u,\cdot).\label{n4:end:conformal-J}
\end{align}
Consequently,
\begin{equation}\label{n4:end:conformal-dec}
 u^2\bigl(\mu_u-|J_u|_{g_u}\bigr)
 \ge \mu-|J|_g+3u^{-1}
          \bigl(-\Delta_g u-|K|_g|\dd u|_g\bigr).
\end{equation}
For any consistently oriented hypersurface,
\begin{equation}\label{n4:end:conformal-expansion}
 \Theta_{K_u}=u^{-1}
       \bigl(\Theta_K+3\partial_\nu\log u\bigr).
\end{equation}
\end{lemma}

\begin{proof}
Use Equation~\eqref{eq:end:conformal} with $n=4$, $k=3$, $p=2$,
and $s=\log u$. The constraint normalization is unchanged. At zero
vector argument its first identity gives
$u^2\mu_u=\mu-3\Delta\log u-3|d\log u|^2
=\mu-3u^{-1}\Delta u$.
The covector formula proved with that identity gives
$uJ_u=J+3K(\nabla\log u,\cdot)$; taking its $g_u$ norm adds
another factor $u^{-1}$. The triangle inequality then gives
Equation~\eqref{n4:end:conformal-dec}, equivalently the $U=u$
specialization of Equation~\eqref{eq:end:conformal-U}.
The boundary formula reads
$\Theta_{K_u}=u^{-1}(\Theta_K+3\partial_\nu\log u)$,
with the same chosen normal. The identities require smooth fields
and $u>0$ only, so this application does not strengthen the original
$O_2/O_1$ decay hypotheses.
\end{proof}

\begin{lemma}[Positive shell]\label{n4:end:positive-shell}
For any one-ended data in Theorem~\ref{n4:thm:weak} and every \(\lambda>0\),
there is a conformal change preserving all its hypotheses, with
\[
 E_\lambda=E+2\lambda,\qquad P_\lambda=P,
 \qquad A_*(g_\lambda)\ge A_*(g).
\]
The conformal factor is constant near the compact boundary and is at
least one everywhere.
\end{lemma}

\begin{proof}
Choose \(0<\xi<\min(q,1)\). For
\(T(r)=r^{-2}-r^{-2-\xi}\), at sufficiently large radius,
\[
 T'<0,\qquad
 -\Delta_gT-|K|_g|\dd T|_g
 =\xi(2+\xi)r^{-4-\xi}+O(r^{-4-q})>0.
\]
Take a nondecreasing smooth switch \(\beta\) from zero to one in this
region, constant on neighborhoods of its endpoints, and set
\[
 G(r)=-\int_r^\infty\beta(s)T'(s)\,\dd s.
\]
Extend \(G\) constantly throughout the interior. Then \(G\ge0\), it
equals \(T\) sufficiently far out, and
\[
 -\Delta_gG-|K|_g|\dd G|_g
 =\beta(-\Delta_gT-|K|_g|\dd T|_g)
       -\beta'T'|\dd r|_g^2\ge0.
\]
Use \(u=1+\lambda G\) in Lemma~\ref{n4:lem:conformal}. The dominant
energy condition is preserved, as is the boundary expansion sign.
The new densities are integrable: outside a compact set the additional
terms in \eqref{n4:end:conformal-mu}--\eqref{n4:end:conformal-J} are bounded
by constants, depending on \(\lambda\), times
\(r^{-4-\xi}+r^{-4-q}\). The new falloff exponent can be taken to be
\(\min(q,2)>1\).

For completeness, the leading change in the metric flux numerator is
\(-6\partial_i(\lambda G)\). Products involving
\(g-\delta\), \(u-1\), or their first derivatives have vanishing
limiting flux. Hence \eqref{n4:intro:energy} gives
\[
 E_\lambda-E=-\frac{\lambda}{\omega_3}
       \lim_{r\to\infty}\int_{S_r}\partial_rG\,\dd A_\delta
       =2\lambda.
\]
The transformed momentum tensor is exactly \(u(K-\tau g)\). Its
additional flux is \(O(r^{-q})\), so \(P_\lambda=P\) by
\eqref{n4:intro:momentum}. Since \(u\ge1\), every cut has at least its
original area, proving the cut comparison.
\end{proof}

The shell will be used only after proving the inequality for
\(E>|P|\). It then excludes all remaining ADM vectors; no preliminary
strict timelikeness hypothesis will remain.

\subsection{Matching the charges before making a rest end}

A rest end must have energy close to $\sqrt{E^2-|P|^2}$, not to $E$.
We therefore first attach a vacuum reference plane carrying exactly
$(E,P)$, and only then bend that plane to a static slice. Compact
linear corrections make the interpolation deficit small; a positive
conformal repair restores DEC at vanishing mass cost.

\subsection*{The reference plane}

\begin{lemma}[Explicit reference-plane charges]\label{n4:end:boosted-plane}
For any \(E>|P|\), put \(m=(E^2-|P|^2)^{1/2}\). A spacelike plane
in the Schwarzschild--Tangherlini spacetime of mass \(m\), sufficiently
far from its central region, has induced vacuum data \((g_*,K_*)\)
with charges exactly \((E,P)\), and
\[
 g_*-\delta=O_k(r^{-2}),\qquad K_*=O_k(r^{-3})
 \quad\text{for every integer }k\ge0.
\]
Here \(K_*\) is one half the metric rate in the future normal direction.
\end{lemma}

\begin{proof}
Use Lemma~\ref{lem:end:boost} with $n=4$, $k=3$, $p=2$, and
$M=m$. Its exact isotropic spacetime metric is
\begin{equation}\label{n4:end:tangherlini}
 \mathbf g=-\left(\frac{1-m/(2|y|^2)}{1+m/(2|y|^2)}\right)^2\mathrm db^2
       +\left(1+\frac{m}{2|y|^2}\right)^2\sum_{i=1}^4(\mathrm dy^i)^2.
\end{equation}
Rotate the first axis into the direction of $P$ and put $v=|P|/E$,
$\gamma=(1-v^2)^{-1/2}$. The Lorentz plane in that lemma has
\(E_*=m\gamma=E\) and \(P_{*1}=m\gamma v=|P|\), with all transverse
momenta zero. These are the present factors $1/(6\omega_3)$ and
$1/(3\omega_3)$: in particular no mass is rescaled. The leading energy
flux is $6m\gamma^2 x_i\rho^{-4}$, the leading momentum tensor is
the $k=3$ instance of \eqref{eq:end:boost-momentum-tensor}, and the
angular integral in \eqref{eq:end:ellipsoid-integral} is
$\omega_3/\gamma$. Thus the sign and normalization are the same
ones used here. The second form uses the same future-normal convention.
Finally the metric and tensor remainders in that lemma are
$O_d(r^{-4})$ and $O_d(r^{-5})$, with leading orders $r^{-2}$ and
$r^{-3}$, through every fixed derivative order. Gauss--Codazzi in the
vacuum metric \eqref{n4:end:tangherlini} proves that the reference data
are vacuum.
\end{proof}

\subsection*{Compact inverses for the linear constraints}

Localized constraint deformation and finite-dimensional adjoint
obstructions have a substantial history
\cite{CorvinoSchoen2006,ChruscielDelay2003}. The compact operators
below are used to reduce, rather than silently solve away, the
interpolation defect. The subsequent DEC repair and comparison of
every admissible cut are separate conclusions of this section.

On Euclidean four-space write
\begin{equation}\label{n4:end:linear-operators}
 \mathcal L b=\partial_i\partial_j
                  (b_{ij}-(\tr_\delta b)\delta_{ij}),\qquad
 (\mathcal D k)_i=\partial_j
                  (k_{ij}-(\tr_\delta k)\delta_{ij}).
\end{equation}
These are the linearizations at \((\delta,0)\) of \(2\mu\) and \(J\).
The next lemma identifies every compatibility condition needed for a
compact correction.

\begin{lemma}[Compact linear constraint corrections]\label{n4:end:compact-inverses}
Fix a compact subannulus \(U_0\) of a bounded connected Euclidean annulus
\(U\). If \(f\) and \(F\) are smooth and supported in \(U_0\), and
\begin{equation}\label{n4:end:compatibility}
 \begin{aligned}
 \int f\phi\,\dd z&=0 &&\text{for every affine }\phi,\\
 \int F_iY^i\,\dd z&=0 &&\text{for every Euclidean Killing field }Y,
 \end{aligned}
\end{equation}
there are symmetric smooth tensors \(b,k\), supported in a fixed compact
subset of \(U\), with \(\mathcal Lb=f\), \(\mathcal Dk=F\).
For every integer \(j\ge0\) and every \(1<p_1<\infty\), they can be
chosen linearly in the sources so that
\begin{equation}\label{n4:end:inverse-bounds}
 \|b\|_{W^{j+2,p_1}}\le C\|f\|_{W^{j,p_1}},\qquad
 \|k\|_{W^{j+1,p_1}}\le C\|F\|_{W^{j,p_1}}.
\end{equation}
The constant depends only on \(U_0,U,j,p_1\) and fixed localization
choices. The moment conditions are also necessary for compactly
supported solutions.
\end{lemma}

\begin{proof}
Lemma~\ref{lem:end:compact-inverses} applies with $n=4$ on the given
connected annulus. Its scalar operator $\mathsf S$ is $\mathcal L$,
and its vector operator $\mathsf D$ is $\mathcal D$. Its affine and
Euclidean Killing moments are exactly \eqref{n4:end:compatibility}.
The inverse trace reversal specializes to
\begin{equation}\label{n4:end:trace-inverse}
 b=S-\tfrac13(\operatorname{tr}_\delta S)\delta,
\end{equation}
so $b-(\operatorname{tr}_\delta b)\delta=S$.
The scalar construction gains two Sobolev derivatives. The symmetric
divergence construction first gains one derivative, uses a further
divergence inverse to correct the skew part, and differentiates that
correction once; its net gain is one. This proves
\eqref{n4:end:inverse-bounds}, with support in a fixed larger compact
subset and linear dependence on the sources. No quantitative derivative
bound for the initial data is being inferred here; the commutator
calculation below establishes the needed source norms from $O_2/O_1$.

For necessity, integration by parts pairs $\mathcal Lb$ with
$\operatorname{Hess}\phi-(\Delta\phi)\delta$, which vanishes for
affine $\phi$, and pairs $\mathcal Dk$ with the symmetric gradient
of a Killing field after trace reversal. The latter also vanishes.
Thus all and only the stated moments obstruct these compact inverses.
\end{proof}

\subsection{Rest-end replacement and comparison of all cuts}

\begin{proposition}[Rest-end replacement]\label{n4:prop:rest-end}
Let \((M,g,K)\) satisfy the one-ended hypotheses of
Theorem~\ref{n4:thm:weak}, and suppose \(E>|P|\). Put
\(m=(E^2-|P|^2)^{1/2}\). There is a sequence of smooth data
\((g_R,K_R)\) on the same manifold, complete with the same compact
boundary, satisfying the dominant energy condition and the same future
boundary inequality, such that:
\begin{enumerate}[label=\textup{(\roman*)}]
 \item \(K_R\) is compactly supported, and the ultimate end is an exact
 static Schwarzschild--Tangherlini spatial end;
 \item its energy is \(m_R=m+o(1)\), its momentum is zero, and both
 constraint densities are integrable;
 \item \(A_*(g_R)\ge(1-o(1))A_*(g)\).
\end{enumerate}
Every error here tends to zero as \(R\to\infty\) for the fixed original
data and fixed timelike vector. No uniformity as \(E-|P|\) tends to zero
is asserted or needed.
\end{proposition}

\begin{proof}
Use Lemma~\ref{n4:end:boosted-plane} for a reference end with exactly the
original charges. We first join the data to this plane on
\(R<r<4R\), controlling the full pointwise constraint deficit.

\paragraph{The scaled commutators.}
Fix \(1<p<\min(2,q)\), write \(x=Rz\), and on
\(1<|z|<4\) use the scaled component fields
\[
 g(Rz),\qquad k(z)=RK(Rz),
\]
and their reference counterparts. This corresponds to rescaling lengths
by \(R^{-1}\), so both constraint densities scale by \(R^2\).
The metric deviations and tensors have size \(O(R^{-p})\) in
\(C^2\) and \(C^1\), respectively. Expanding the constraint map about
\((\delta,0)\) gives its linear part
\((\mathcal L,\mathcal D)\) from \eqref{n4:end:linear-operators}, with
remainder bounded in \(C^0\) by \(CR^{-2p}\). This follows directly
from the scalar curvature formula: the nonlinear metric terms are
bounded by \(C(|b||D^2b|+|Db|^2)\); the other scalar terms are quadratic
in \(k\). The momentum remainder is bounded by
\(C(|b||Dk|+|Db||k|)\).

Choose a fixed smooth radial \(\psi\), equal to one near \(|z|=1\)
and zero near \(|z|=4\), with \(0\le\psi\le1\). Blend the metric
components and tensors by \(\psi\). If \(b\) is the difference of
the two scaled metric perturbations and \(k\) the difference of the
two scaled tensors, the errors of the linear constraints relative to
the blended sources are
\(f_R=[\mathcal L,\psi]b\) and
\(F_R=[\mathcal D,\psi]k\). For \(T_{ij}=b_{ij}-(\tr_\delta b)\delta_{ij}\),
\begin{align}
 [\mathcal L,\psi]b
   &=(\partial_i\partial_j\psi)T_{ij}
       +2(\partial_i\psi)\partial_jT_{ij},\label{n4:end:scalar-commutator}\\
 ([\mathcal D,\psi]k)_i
   &=(\partial_j\psi)(k_{ij}-(\tr_\delta k)\delta_{ij}).
       \label{n4:end:vector-commutator}
\end{align}
These expressions are supported in a fixed compact subannulus, and
their \(C^1\) norms are \(O(R^{-p})\). In particular, the first bound
uses only two metric derivatives, and the second only one tensor
derivative.

\paragraph{All compatibility moments are small.}
We prove the sharper bound \(o(R^{-2})\) for every moment in
\eqref{n4:end:compatibility}. Work momentarily in the physical coordinates,
with
\[
 h=g-g_*,\quad
 T=K-(\tr_\delta K)\delta-K_*+(\tr_\delta K_*)\delta,
 \quad f=\mathcal Lh,\quad F_i=\partial_jT_{ij}.
\]
Both \(f\) and \(F\) belong to \(L^1\) on the end. Indeed, the
difference between the physical constraints and their Euclidean linear
parts is \(O(r^{-2-2\min(q,2)})\), which is integrable in dimension
four. The physical densities are integrable by hypothesis and by
reference vacuum, and Euclidean and physical measures and norms are
uniformly comparable there.

An elementary consequence of integrability is
\begin{equation}\label{n4:end:weak-first-moment}
 \int_{r_0<r<4R}r(|f|+|F|)\,\dd x=o(R).
\end{equation}
To see this, split at a fixed radius \(L\). The integral up to \(L\),
divided by \(R\), tends to zero. The remaining integral divided by
\(R\) is at most four times the \(L^1\) tail beyond \(L\), which can
be made arbitrarily small. A finite first moment is not required.

For an affine function \(\phi\), define the flux primitive
\[
 Q_{\phi,i}
 =\phi(\partial_jh_{ij}-\partial_i h_{jj})
       -(\partial_j\phi)h_{ij}+(\partial_i\phi)h_{jj}.
\]
Its divergence is \(\phi f\). For a Killing field \(Y\), symmetry
of \(T\) gives \(\partial_j(Y^iT_{ij})=Y^iF_i\).
Let \(\mathfrak F_\phi(r)\) and \(\mathfrak G_Y(r)\) be the corresponding
outward Euclidean sphere fluxes. For constant tests, matching the ADM
charges gives
\[
 \mathfrak F_1(r)=o(1),\qquad
 \mathfrak G_{e_i}(r)=o(1).
\]
In the momentum flux the replacement of \(\tr_gK\,g\) by
\((\tr_\delta K)\delta\) costs at most
\(O(r^{2-2\min(q,2)})=o(1)\). For homogeneous linear scalar tests and
rotational fields, integration from a fixed sphere and
\eqref{n4:end:weak-first-moment} instead give
\[
 \mathfrak F_\phi(r)=o(r),\qquad \mathfrak G_Y(r)=o(r).
\]

Put \(\psi_R(x)=\psi(x/R)\), and let \(A_R=\{R<r<4R\}\).
Integration by parts, with \(\psi_R=1\) near the inner sphere and zero
near the outer sphere, gives the exact formulas
\begin{align}
 \int_{A_R}\phi[\mathcal L,\psi_R]h\,\dd x
  &=-\mathfrak F_\phi(R)-\int_{A_R}\psi_R\phi f\,\dd x,
       \label{n4:end:scalar-moment}\\
 \int_{A_R}Y^i([\mathcal D,\psi_R](K-K_*))_i\,\dd x
  &=-\mathfrak G_Y(R)-\int_{A_R}\psi_RY^iF_i\,\dd x.
       \label{n4:end:vector-moment}
\end{align}
In scaled coordinates, constant moments have the factor \(R^{-2}\):
the source scales by \(R^2\) and volume by \(R^{-4}\).
A linear test in \(z\) contributes the additional factor \(R^{-1}\).
Thus \eqref{n4:end:scalar-moment}--\eqref{n4:end:vector-moment} and the preceding
estimates show
\begin{equation}\label{n4:end:small-moments}
 \int f_R\phi\,\dd z=o(R^{-2}),\qquad
 \int (F_R)_iY^i\,\dd z=o(R^{-2})
\end{equation}
for each element of a fixed basis of the affine functions or Killing
fields.

\paragraph{Correction and the physical deficit.}
Choose a nonnegative smooth bump \(\zeta\), supported in a fixed ball
within the annulus and positive on a smaller ball. For a basis of the
affine tests, its Gram matrix
\(\int\zeta\phi_a\phi_b\) is positive definite. For a basis of
Killing fields, the matrix \(\int\zeta\,Y_a\cdot Y_b\) is likewise
positive definite: a nonzero affine Killing field cannot vanish on an
open ball. Subtract the corresponding bump combinations from
\(f_R,F_R\) to erase their moments. By \eqref{n4:end:small-moments},
the coefficients, and every fixed smooth norm of these bump corrections,
are \(o(R^{-2})\).

Apply Lemma~\ref{n4:end:compact-inverses} to the negatives of the
moment-free commutators and add the resulting tensors to the blend.
With \(j=1\) and \(p_1>4\), Sobolev embedding and
\eqref{n4:end:inverse-bounds} bound the metric correction in \(C^2\)
and the tensor correction in \(C^1\) by \(CR^{-p}\).
Their supports stay in a fixed enlarged compact subannulus.
The new metric is positive for sufficiently large \(R\) and agrees
exactly with the original data inside \(R\) and the plane outside
\(4R\).

Let \((\widetilde g_R,\widetilde K_R)\) denote these physical data.
Their scaled constraints equal \(\psi\) times the original scaled
constraints, plus an error
\(o(R^{-2})+O(R^{-2p})\); the reference constraints vanish.
The change in the momentum norm costs another \(O(R^{-2p})\), because
the metric change is \(O(R^{-p})\) and the scaled momentum density is
\(O(R^{-p})\). Since \(\psi\ge0\), the original dominant energy
condition implies, after undoing the scaling,
\begin{equation}\label{n4:end:physical-deficit}
 \widetilde\mu_R-|\widetilde J_R|_{\widetilde g_R}
       \ge-\eta_RR^{-4},\qquad \eta_R\longrightarrow0.
\end{equation}
Indeed the remaining physical error is
\(R^{-2}[o(R^{-2})+O(R^{-2p})]=o(R^{-4})\), since \(p>1\).
The possible deficit is confined to a fixed closed subannulus on the
scale \(R\). Replace \(\eta_R\) by a positive majorant tending to zero
if necessary.

\paragraph{Bending the plane to a static slice.}
Only the exact reference region is used for this step. In its spatial
coordinates \(y\), write the plane as \(b=v\cdot y\). Choose a smooth
switch \(\psi_1\) from one to zero on \([0,1]\), constant outside that
interval, and replace the plane farther out by the graph
\begin{equation}\label{n4:end:vacuum-bend}
 b=(v\cdot y)\psi_1\left(L^{-1}\log(|y|/R_b)\right).
\end{equation}
Here \(R_b\) is a sufficiently large fixed multiple of \(R\), so the
graph initially lies wholly beyond the gluing annulus. For fixed
\(|v|<1\), choose the fixed \(L\) so large that
\[
 |\grad_y b|
 \le |v|\bigl(1+\|\psi_1'\|_\infty/L\bigr)<1.
\]
The Minkowski graph is uniformly spacelike. Since the perturbation
\eqref{n4:end:tangherlini} is \(O(r^{-2})\), it remains uniformly spacelike
in that metric for large \(R\). It agrees with the plane on an open
inner region and with \(b=0\) on an open outer region. These induced
data are vacuum everywhere in the bend, so it adds no deficit.

In the fixed \(x\) chart, use \(y=Ax\) where
\(A=(I-v\otimes v)^{-1/2}\), with the chosen rotation. The derivatives
of \eqref{n4:end:vacuum-bend} give \(D^2b=O(r^{-1})\), and the induced
data, now denoted \((g_R^0,K_R^0)\), satisfy on all transition regions
\begin{equation}\label{n4:end:transition-geometry}
 c\delta\le g_R^0\le C\delta,\qquad
 |\partial g_R^0|+|K_R^0|\le C/r.
\end{equation}
Here \(c,C>0\) may depend on the fixed boost and switches, but not
on large \(R\). Outside a fixed multiple of \(R\), the data are static
in the \(y\) chart and have tensor zero.

For the repair choose a smooth radius \(\mathfrak r\), equal to
\(|x|\) through the gluing annulus and through the bend, and equal to
\(|y|\) farther out in the static region. It can be chosen with
\begin{equation}\label{n4:end:radius-bounds}
 \mathfrak r\asymp r,\qquad
 c\le|\dd\mathfrak r|_{g_R^0}\le C,
 \qquad |\Delta_{g_R^0}\mathfrak r|\le C/r.
\end{equation}
Here is an explicit way to make the interpolation. In the static region
put \(a(n)=|A^{-1}n|\), \(n=y/|y|\), and interpolate
\[
 \log\mathfrak r=\log|y|+(1-\chi(\log(|y|/R_c)))\log a(n).
\]
Take \(R_c\) beyond the bend and make \(\chi\) change from zero to one
over a sufficiently long fixed logarithmic interval. Its slope can be
made so small that the derivative of \(\log\mathfrak r\) with respect
to \(\log|y|\) lies between \(1/2\) and \(3/2\). Angular derivatives
are bounded, second spatial derivatives are \(O(r^{-1})\), and
\eqref{n4:end:transition-geometry} proves \eqref{n4:end:radius-bounds}.
All transition intervals are contained between fixed multiples of \(R\).

\paragraph{A conformal repair with vanishing mass cost.}
Write \(s=\mathfrak r/R\). We construct a nonnegative function
\(F_R(s)\), constant on the inside, and use
\[
 u_R=1+\varepsilon_RR^{-2}F_R(s),\qquad
 \varepsilon_R=C_0\eta_R.
\]
Choose a fixed interval \([a,b_1]\) in the \(s\) variable that starts
before the possible deficit, contains all transition regions, and ends
where the radius is \(|y|\) and the data are static. This interval lies
in the end for all large \(R\). Put \(z_R=-F_R'\) and prescribe
\begin{equation}\label{n4:end:repair-ode}
 z_R(a)=0,\qquad z_R'=C_1z_R+\omega_1,
\end{equation}
where \(\omega_1\ge0\) is smooth, vanishes on a neighborhood of \(a\),
and is at least one on a neighborhood of the possible deficit interval.
Choose it to vanish near \(b_1\). The solution is nonnegative and its
norms on this fixed interval are bounded independently of large \(R\).

Differentiating \(u_R\) and using \eqref{n4:end:transition-geometry} and
\eqref{n4:end:radius-bounds} gives fixed constants \(c',C'>0\) such that
\begin{equation}\label{n4:end:repair-estimate}
 -\Delta_{g_R^0}u_R-|K_R^0||\dd u_R|_{g_R^0}
 \ge\varepsilon_RR^{-4}(c'z_R'-C'z_R)
 \quad\text{on }a\le s\le b_1.
\end{equation}
In fact the first positive term is
\(\varepsilon_RR^{-4}z_R'|\dd\mathfrak r|^2\); the other terms are
\(\varepsilon_RR^{-3}z_R\Delta\mathfrak r\) and the negative tensor
term, both bounded below by \(-C\varepsilon_RR^{-4}z_R\) here.
Choose \(C_1>C'/c'\). Equation \eqref{n4:end:repair-ode} then makes
\eqref{n4:end:repair-estimate} nonnegative everywhere on this interval
and at least \(c'\varepsilon_RR^{-4}\) on the deficit region.

On the static region put \(c_m=m/2\) and continue \(z_R\) so that
\[
 q_R(s)=s^3\left(1+\frac{c_m}{R^2s^2}\right)^2z_R(s)
\]
is nondecreasing and becomes a positive constant \(q_{R,\infty}\).
This can be done smoothly with a uniform bound on \(q_{R,\infty}\).
Indeed, immediately past \(b_1\), continue \eqref{n4:end:repair-ode};
both its derivative and that of the positive prefactor make \(q_R'\)
nonnegative for large \(R\). Multiply this nonnegative derivative by
a smooth switch that equals one initially and zero after one more fixed
interval, and integrate. This preserves all matching derivatives and
keeps \(q_R\) nondecreasing. The radial Laplacian of the static metric
then shows \(-\Delta_{g_R^0}u_R\ge0\) there, since its radial divergence
is proportional to \(-q_R'\); also \(K_R^0=0\).

Set \(F_R(s)=\int_s^\infty z_R(t)\,\dd t\), extending it constantly
before \(a\). Its norm is bounded uniformly: the finite intervals have
bounded length and data, while on the tail \(z_R=O(s^{-3})\).
Thus \(1\le u_R\le2\) for large \(R\).
Choose the fixed \(C_0\) large enough that
\(3u_R^{-1}c'\varepsilon_R\ge\eta_R\). Lemma~\ref{n4:lem:conformal},
\eqref{n4:end:physical-deficit}, and \eqref{n4:end:repair-estimate} now prove
the dominant energy condition everywhere for
\[
 g_R=u_R^2g_R^0,\qquad K_R=u_RK_R^0.
\]
The repair is constant in a neighborhood of the original boundary, and
therefore preserves its expansion sign.

On the ultimate tail the integral for \(F_R\) is explicit:
\[
 u_R=1+\frac{a_R}{|y|^2+c_m},\qquad
 a_R=\tfrac12\varepsilon_Rq_{R,\infty}=O(\varepsilon_R)=o(1).
\]
Consequently
\begin{equation}\label{n4:end:static-repaired-metric}
 g_R=\left(1+\frac{c_m+a_R}{|y|^2}\right)^2\delta,
 \qquad K_R=0
 \quad\text{on that tail}.
\end{equation}
Its energy is \(m+2a_R=m+o(1)\), with zero momentum, by the direct
static instance of Lemma~\ref{n4:end:boosted-plane}. The densities vanish
on this tail and are smooth on the remaining compact part, so they are
integrable. The data are complete with boundary: a fixed compact core
is smooth, and the end has a uniform positive metric lower bound for
each \(R\).

\paragraph{Comparison of every enclosing cut.}
The final metric dominates the original \(g\) for \(r\le R\), because
the intermediate data there are the original ones and \(u_R\ge1\).
For \(r\ge R\), \eqref{n4:end:transition-geometry} and the static tail
give \(g_R\ge c\delta\) in the original \(x\) chart with a fixed
\(c>0\). Choose \(0<\kappa<\min(1,\sqrt c/2)\).
Let \(L_R=R^{1/2}\), and choose a smooth strictly increasing radial map
\(\rho_R\) equal to \(r\) for \(r\le L_R\), with derivative decreasing
from one to \(\kappa\) on \([L_R,2L_R]\), and derivative \(\kappa\)
afterward. It defines a diffeomorphism
\[
 \Phi_R(x)=\rho_R(|x|)\,\frac{x}{|x|}
\]
on the end, extended by the identity on the core. It is proper because
\(\rho_R(r)\to\infty\), preserves the selected end, and fixes the
boundary. Its radial and tangential Euclidean dilations are
\(\rho_R'\) and \(\rho_R/r\). They are at most one everywhere, and
for \(r\ge R\) are at most \(\kappa+O(R^{-1/2})\).

On the nonidentity region with \(r\le R\), both \(r\) and
\(\rho_R(r)\) tend uniformly to infinity; the original metric is
therefore \((1+o(1))\delta\) at both points. The dilation bound and
\(g_R\ge g\) show
\(\Phi_R^*g\le(1+o(1))g_R\) there. For \(r\ge R\), the choice
of \(\kappa\), the same original-end decay, and \(g_R\ge c\delta\)
give this inequality as well. It is immediate on the identity core.
We have thus proved the global tensor comparison
\begin{equation}\label{n4:end:compression-metric}
 \Phi_R^*g\le(1+\epsilon_R)g_R,\qquad\epsilon_R\longrightarrow0.
\end{equation}
Let \(\Gamma\) be the full boundary of any admissible outer domain
for \(g_R\). Its image under this proper diffeomorphism is the full
boundary of an admissible original outer domain, with all coincident
frontier retained. Equation~\eqref{n4:end:compression-metric} gives
\[
 |\Phi_R(\Gamma)|_g
 \le(1+\epsilon_R)^{3/2}|\Gamma|_{g_R}.
\]
Taking infima gives
\[
 A_*(g)\le(1+\epsilon_R)^{3/2}A_*(g_R),
\]
which proves the last assertion without a compactness assumption on
minimizing sequences.
\end{proof}

\subsection{Strict trapping and a positive energy margin}
\label{n4:sec:strictification}

The rest-end replacement preserves weak trapping and the dominant energy
condition.  The prepared construction requires strict inequalities.  We
now obtain them by a conformal change whose size tends uniformly to zero
and whose change in ADM energy can be computed exactly.  All constants
in this step belong to one fixed rest-end metric.  In particular, when
that metric depends on the replacement radius \(R\), we keep \(R\)
fixed throughout this construction.

The drift-Poisson equation used below is related to Jaracz's conformal
strictification~\cite{Jaracz2025StrictDEC}.  The proof here supplies the
four-dimensional equation, its nonzero inner Neumann condition, and the
charge and full-cut estimates needed for the numerical argument.

\begin{lemma}[Strictification of a fixed rest end]
\label{n4:lem:strictification}
Let \((M,g,K)\) be a smooth connected four-dimensional exterior,
complete with its nonempty compact smooth boundary included, with one
Euclidean coordinate end and compact complement.  Suppose
\(\mu\ge |J|_g\) and \(\Theta_K\le0\) on every boundary component,
where the boundary normal \(\nu\) points into \(M\).  Assume that
\(K\) has compact support and that, outside a compact set,
\[
 g=\left(1+\frac{c}{r^2}\right)^2\delta
\]
in the end coordinates, for a constant \(c\).
Fix \(0<\delta_1<1\), and extend the coordinate radius to a positive
smooth function \(r\) on \(M\).

There is a nonnegative smooth function \(\phi\), with
\(\phi=O_k(r^{-2})\) for every fixed derivative order \(k\), such
that, for all sufficiently small \(s>0\), the data
\[
 (g_s,K_s)=(e^{2s\phi}g,e^{s\phi}K)
\]
satisfy
\begin{equation}\label{n4:eq:strictification-output}
 \Theta_{K_s}<0,\qquad
 \mu_s-|J_s|_{g_s}\ge c_s r^{-4-\delta_1}
 \quad\hbox{on }M,\qquad c_s>0.
\end{equation}
They are complete, have compactly supported second fundamental form,
integrable constraint densities, zero ADM momentum and finite ADM
energy.  Their end satisfies
\begin{equation}\label{n4:eq:strictification-decay}
 g_s-\delta=O_k(r^{-2})\quad\hbox{for every fixed }k,
 \qquad R_{g_s}=O(r^{-4-\delta_1}).
\end{equation}
For the full enclosing-cut infimum on this fixed exterior,
\begin{equation}\label{n4:eq:strictification-limits}
 E_{g_s}\longrightarrow E_g,
 \qquad A_*(g_s)\longrightarrow A_*(g)
 \quad\hbox{as }s\downarrow0.
\end{equation}
\end{lemma}

\begin{proof}
The generic strictification lemma, Lemma~\ref{lem:end:strict}, applies
to this fixed rest exterior. Indeed the exact static tail has
$g-\delta=O_d(r^{-2})$ for every $d$, and $K$ is compactly supported.
Choose any $1<q_*<2$ to express its $O_6/O_5$ bounds, set $n=4$,
$k=3$, and take $\beta=\delta_1$. The condition
$0<\delta_1<\min(1,q_*)$ is exactly the one required by that lemma.
Its smooth nonnegative majorant can be chosen compactly supported:
write it as $b_0\ge|K|_g$. The complete mixed-boundary continuation,
uniform core and tail estimates, exhaustion, and differentiated-decay
proof of that lemma therefore give a positive smooth solution of
\begin{equation}\label{n4:eq:strictification-pde}
 -\Delta_g\phi=b_0\sqrt{|\mathrm d\phi|_g^2+r^{-6}}+r^{-4-\delta_1},
 \qquad \partial_\nu\phi=-1,\qquad \phi\longrightarrow0.
\end{equation}
The exponent $-6$ is $-2k$. The all-order clause of the same lemma
gives
\begin{equation}\label{n4:eq:strictification-phi-decay}
 |\partial^j\phi|\le C_jr^{-2-j}\qquad(j\ge0).
\end{equation}
This application occurs after rest replacement, so it uses no higher
derivatives of the original weakly asymptotically flat end.

We now verify the geometric conclusions. Put
\[F=b_0\sqrt{|\mathrm d\phi|^2+r^{-6}}+r^{-4-\delta_1}.\]
The inner outward normal is $-\nu$, and the finite-flux conclusion of the lemma is
\begin{equation}\label{n4:eq:strictification-flux}
 L_\phi=\lim_{T\to\infty}\int_{S_T}\partial_r\phi\,\mathrm dA_\delta
       =-|\partial M|_g-\int_MF\,\mathrm dV_g<0.
\end{equation}
The geometric and Euclidean fluxes differ by $O(T^{-2})$ on the static
tail, by \eqref{n4:eq:strictification-phi-decay}.

Apply Lemma~\ref{n4:lem:conformal} with $u=e^{s\phi}$. It gives
\[
 e^{2s\phi}(\mu_s-|J_s|_{g_s})
 \ge3s\bigl(r^{-4-\delta_1}-s|\mathrm d\phi|^2\bigr).
\]
The number $B=\sup_M r^{4+\delta_1}|\mathrm d\phi|^2$ is finite:
smoothness controls the compact core and the tail is $O(r^{-2+\delta_1})$.
Choose $sB\le1/2$. This proves the global margin with
$c_s=(3s/2)e^{-2s\|\phi\|_\infty}>0$. On each boundary component,
\[
 \Theta_{K_s}=e^{-s\phi}(\Theta_K+3s\partial_\nu\phi)
            =e^{-s\phi}(\Theta_K-3s)<0.
\]
Furthermore
\begin{equation}\label{n4:eq:strictification-metric-comparison}
 g\le g_s\le e^{2s\|\phi\|_\infty}g.
\end{equation}
Thus completeness with the boundary included is preserved, and the
tensor remains compactly supported. On the tail, $R_g=0$, $K_s=0$ and
\[
 R_{g_s}=6e^{-2s\phi}
       \bigl(sr^{-4-\delta_1}-s^2|\mathrm d\phi|^2\bigr)
       =O(r^{-4-\delta_1}).
\]
This proves the scalar decay and integrability there; the rest is smooth
and compact. All differentiated metric bounds follow from
\eqref{n4:eq:strictification-phi-decay}, and momentum is zero because
the tensor vanishes on a distant tail.

The leading metric change is $2s\phi\delta+O_1(r^{-4})$, whose energy
flux numerator is $-6s\partial_i\phi+O(r^{-5})$. The latter remainder
has vanishing sphere flux. Division by $6\omega_3$ therefore gives
\begin{equation}\label{n4:eq:strictification-energy}
 E_{g_s}-E_g=-\frac{s}{\omega_3}L_\phi\longrightarrow0.
\end{equation}
Every full cut has its three-dimensional area element multiplied by
$e^{3s\phi}$; the cut class itself is unchanged. Taking infima yields
\begin{equation}\label{n4:eq:strictification-area}
 A_*(g)\le A_*(g_s)\le e^{3s\|\phi\|_\infty}A_*(g).
\end{equation}
The infimum is finite, since a distant coordinate sphere is admissible.
Its convergence as $s\downarrow0$ follows without an attained minimizer.
Together these estimates prove every asserted output at this fixed rest
datum. If it is $(g_R,K_R)$, every constant above may depend on $R$;
strictification is removed with $R$ fixed.
\end{proof}

\subsection{Applying the prepared theorem and returning to the original end}

The proper radial comparison protects the original enclosing area during
rest-end replacement. Strictification then changes the cut infimum by a
factor tending to one at each fixed replacement radius. We apply the
prepared inequality after both constructions, remove strictification at
that fixed radius, and finally send the replacement radius to infinity.

\begin{proposition}[Transfer from prepared data]\label{n4:end:transfer}
Theorem~\ref{thm:numerical:prepared}, applied in dimension four,
implies Theorem~\ref{n4:thm:weak}.
\end{proposition}

\begin{proof}
First let the original exterior have one end and $E>|P|$.
Set $m=(E^2-|P|^2)^{1/2}>0$.
Proposition~\ref{n4:prop:rest-end} gives data $(g_R,K_R)$ with
zero momentum and energy $m_R\to m$. Its full-cut comparison is
\begin{equation}\label{n4:end:one-sided-area}
 A_*(g_R)\ge(1+\epsilon_R)^{-3/2}A_*(g),
 \qquad \epsilon_R\longrightarrow0.
\end{equation}
Fix one sufficiently large replacement radius $R$. For any fixed
$0<\delta_1<1$, Lemma~\ref{n4:lem:strictification} makes
$(g_{R,s},K_{R,s})$ a prepared exterior for all sufficiently small
$s>0$. Indeed it is the same smooth connected orientable one-ended
exterior, complete with its compact boundary included and with compact
complement of its coordinate end. It has finite ADM energy and a
compactly supported tensor, so its ADM momentum is zero. On the end,
$g_{R,s}-\delta=O_d(r^{-2})$ for every fixed derivative order $d$,
its scalar curvature is
$O(r^{-4-\delta_1})$, and its constraint densities are integrable.
The boundary is strictly future trapped, and the global margin is
\[
 \mu_{R,s}-|J_{R,s}|_{g_{R,s}}
 \ge c_{R,s}\rho^{-4-\delta_1},\qquad c_{R,s}>0,
\]
where $\rho$ is a positive smooth extension of the rest-end radius.
These are the geometric and decay conditions of
Definition~\ref{def:prepared}, with $n=4$ and $\beta=\delta_1$.

The area condition follows along the same comparisons. In one end,
Definition~\ref{n4:def:cuts} is exactly the prepared full-cut class:
the outer domain is connected and closed, its interior lies in the
open exterior, and its entire compact intrinsic boundary is counted,
including every component and all coincidence with the original boundary.
Lemma~\ref{n4:lem:cut-reduction} gives $A_*(g)>0$;
Equation~\eqref{n4:end:one-sided-area} and strictification give
$A_*(g_{R,s})\ge A_*(g_R)>0$.
The constraint and ADM normalizations also agree:
$2\mu=R+(\operatorname{tr}K)^2-|K|^2$,
$J=\operatorname{div}(K-(\operatorname{tr}K)g)$, and the energy and
momentum factors are $1/(6\omega_3)$ and $1/(3\omega_3)$.
No rescaling is made.

Theorem~\ref{thm:numerical:prepared} therefore yields, at every such
fixed pair $(R,s)$,
\begin{equation}\label{n4:end:fixed-prepared-inequality}
 E_{R,s}\ge\frac12
       \left(\frac{A_*(g_{R,s})}{\omega_3}\right)^{2/3}.
\end{equation}
This is the completed numerical theorem: its elliptic exhaustion and
its limits $N\to\infty$ and $\epsilon\downarrow0$ have already been
taken for the fixed prepared datum. None of those estimates is required
to be uniform in $R$ or $s$.

Now let $s\downarrow0$ with $R$ fixed. Strictification gives both
$E_{R,s}\to m_R$ and $A_*(g_{R,s})\to A_*(g_R)$, so
\[
 m_R\ge\frac12
       \left(\frac{A_*(g_R)}{\omega_3}\right)^{2/3}
 \ge\frac12(1+\epsilon_R)^{-1}
       \left(\frac{A_*(g)}{\omega_3}\right)^{2/3}.
\]
Only at this point let $R\to\infty$. The result is
\begin{equation}\label{n4:end:timelike-conclusion}
 \sqrt{E^2-|P|^2}\ge\frac12
       \left(\frac{A_*(g)}{\omega_3}\right)^{2/3}
       \qquad (E>|P|).
\end{equation}
Full convergence of $A_*(g_R)$ has not been used or asserted.

It remains to remove the timelike assumption. If instead $E\le|P|$,
take any sequence $a_j>0$ tending to zero and set
\[
 \lambda_j=\frac{|P|+a_j-E}{2}>0.
\]
Lemma~\ref{n4:end:positive-shell} gives, separately for each $j$,
data of energy $|P|+a_j$, momentum $P$, and enclosing infimum at
least $A_*(g)$. Their charges are future timelike, so
\eqref{n4:end:timelike-conclusion} gives
\[
 \sqrt{(|P|+a_j)^2-|P|^2}
 \ge\frac12\left(\frac{A_*(g)}{\omega_3}\right)^{2/3}>0.
\]
The right side is fixed and positive by
Lemma~\ref{n4:lem:cut-reduction}, while the left side tends to zero.
This contradiction includes zero, negative, and null energy cases.
No preparation geometry is passed to a limit as the boost degenerates:
the timelike numerical result was applied afresh to each shell datum.
Thus $E>|P|$ and \eqref{n4:end:timelike-conclusion} hold for every
one-ended original datum in the theorem.

Finally let the original manifold have finitely many ends and fix
one end $e$. Truncate every other end at a sufficiently large sphere
as in Lemma~\ref{n4:lem:cut-reduction}. The retained manifold is a
complete connected one-ended exterior. Its added boundary components
are strictly future trapped, its selected-end charges are still
$(E_e,P_e)$, and its full-cut infimum is at least $A_e(S)$. Applying
the one-ended result gives
\[
 E_e\ge\sqrt{|P_e|^2+
              \frac14\left(\frac{A_e(S)}{\omega_3}\right)^{4/3}}.
\]
Positivity of $A_e(S)$ was proved before any truncation. This gives
both the asserted inequality and strict future timelikeness at $e$.
\end{proof}

The one-ended cases of Theorems~\ref{thm:weak3} and~\ref{n4:thm:weak}
together prove Theorem~\ref{thm:low}.

\subsection{The general maximal DEC energy inequality}

The maximal special case requires no assumption that a momentum limit
exists. The following end calculation supplies that limit from the
decay and integrability already imposed on the data.

\begin{lemma}[Existence of the momentum flux]\label{n4:lem:momentum-exists}
On a smooth four-dimensional asymptotically flat end, suppose
$g-\delta=O_2(r^{-q})$ and $K=O_1(r^{-1-q})$ for some $q>1$.
If $|J|_g$ is integrable, where
$J_i=\nabla^j(K_{ij}-(\tr_gK)g_{ij})$, then each ADM momentum
limit in \eqref{n4:intro:momentum} exists and is finite.
\end{lemma}

\begin{proof}
In the given coordinates set
\[
 \pi^g_{ij}=K_{ij}-(\tr_gK)g_{ij},\qquad
 \pi^\delta_{ij}=K_{ij}-(\tr_\delta K)\delta_{ij}.
\]
The differentiated decay assumptions give
\begin{equation}\label{n4:end:momentum-integrability}
 \pi^g-\pi^\delta=O_1(r^{-1-2q}),\qquad
 \partial_j\pi^\delta_{ij}-J_i=O(r^{-2-2q}).
\end{equation}
To verify the second bound, expand the covariant divergence as
\[
 J_i=g^{jk}\bigl(\partial_k\pi^g_{ij}
       -\Gamma^a_{ki}\pi^g_{aj}-\Gamma^a_{kj}\pi^g_{ia}\bigr).
\]
Replacing $g^{jk}$ by $\delta^{jk}$ costs
$O(r^{-q})O(r^{-2-q})$; replacing $\pi^g$ by $\pi^\delta$
costs the derivative of the first bound in
\eqref{n4:end:momentum-integrability}. Each connection term has size
$O(r^{-1-q})O(r^{-1-q})$. These are exactly the asserted errors.

Their Euclidean radial volume bound is $Cr^{1-2q}$, integrable because
$q>1$. The asymptotic metric and Euclidean volume and covector norms
are uniformly comparable, so $\partial_j\pi^\delta_{ij}\in L^1$.
The Euclidean divergence theorem between two large spheres makes
the flux $\int_{S_R}\pi^\delta_{ij}n_\delta^j\,\dd A_\delta$
a Cauchy function of $R$, with finite limit. By the first bound in
\eqref{n4:end:momentum-integrability}, its difference from the flux
of $\pi^g$ is $O(R^{2-2q})\to0$. This is the tensor appearing
in \eqref{n4:intro:momentum}, which proves the lemma.
\end{proof}

\begin{corollary}[Maximal DEC energy inequality]
\label{n4:cor:maximal-dec}\label{n4:ass:energy}
Let $(X^4,g_X,K_X)$ be smooth and connected, with $X$ orientable,
$g_X$ complete as a metric space with its nonempty compact smooth
boundary $S_X$ included, and with exactly one end, which is asymptotically
flat and diffeomorphic to the complement of a closed ball in $\R^4$.
Let $K_X$ be a smooth symmetric covariant two-tensor. Suppose
\[
\begin{gathered}
 \tr_{g_X}K_X=0,\qquad
 \mu_X=\tfrac12(R_{g_X}-|K_X|_{g_X}^2)
       \ge|\divg_{g_X}K_X|_{g_X},\\
 \mu_X,|\divg_{g_X}K_X|_{g_X}\in L^1(X,\dd V_{g_X}).
\end{gathered}
\]
Assume for some $q>1$ that
\[
 g_X-\delta=O_2(r^{-q}),\qquad K_X=O_1(r^{-1-q}),
\]
and that the ADM energy $E_X$ in \eqref{n4:intro:energy} is finite.
On every component of $S_X$, require
$H+\tr_{S_X}K_X\le0$ with the normal into $X$.
For the one-ended enclosing-cut infimum of Definition~\ref{n4:def:cuts},
denoted here by $a_{g_X}(S_X)$, one has
\[
 E_X\ge\frac12\left(\frac{a_{g_X}(S_X)}{2\pi^2}\right)^{2/3}.
\]
The boundary may have any topology and finitely many components.
No symmetry, vacuum, spin, outermostness, outer area-minimization
or momentum-flux assumption is imposed.
\end{corollary}

\begin{proof}
Maximality identifies the momentum density with
$J_X=\divg_{g_X}K_X$. Lemma~\ref{n4:lem:momentum-exists} supplies
the finite ADM momentum. All hypotheses of
Theorem~\ref{n4:thm:weak} therefore hold for this one-ended exterior.
Its invariant-mass inequality in particular gives the displayed
energy bound. The proof of Theorem~\ref{n4:thm:weak} uses
Theorem~\ref{thm:numerical:prepared}, independently of this corollary.
This numerical consequence retains its broad maximal-DEC scope; the
separate direct maximal-vacuum construction has its own hypotheses.
\end{proof}

